\nonstopmode{}
\documentclass[a4paper]{book}
\usepackage[times,inconsolata,hyper]{Rd}
\usepackage{makeidx}
\makeatletter\@ifl@t@r\fmtversion{2018/04/01}{}{\usepackage[utf8]{inputenc}}\makeatother
\usepackage{graphicx}\setkeys{Gin}{width=0.7\textwidth}
\makeindex{}
\begin{document}
\chapter*{}
\begin{center}
{\textbf{\huge Package `ggalt'}}
\par\bigskip{\large \today}
\end{center}
\ifthenelse{\boolean{Rd@use@hyper}}{\hypersetup{pdftitle = {ggalt: Extra Coordinate Systems, 'Geoms', Statistical Transformations, Scales and Fonts for 'ggplot2'}}}{}
\ifthenelse{\boolean{Rd@use@hyper}}{\hypersetup{pdfauthor = {Bob Rudis; Ben Bolker; Jan Schulz; Aditya Kothari; Jonathan Sidi}}}{}
\begin{description}
\raggedright{}
\item[Title]\AsIs{Extra Coordinate Systems, 'Geoms', Statistical Transformations,
Scales and Fonts for 'ggplot2'}
\item[Version]\AsIs{0.6.1}
\item[Maintainer]\AsIs{Bob Rudis }\email{bob@rud.is}\AsIs{}
\item[Description]\AsIs{A compendium of new geometries, coordinate systems, statistical 
transformations, scales and fonts for 'ggplot2', including splines, 1d and 2d densities, 
univariate average shifted histograms, a new map coordinate system based on the 
'PROJ.4'-library along with geom_cartogram() that mimics the original functionality of 
geom_map(), formatters for ``bytes'', a stat_stepribbon() function, increased 'plotly'
compatibility and the 'StateFace' open source font 'ProPublica'. Further new 
functionality includes lollipop charts, dumbbell charts, the ability to encircle
points and coordinate-system-based text annotations.}
\item[License]\AsIs{MIT + file LICENSE}
\item[URL]\AsIs{}\url{https://github.com/hrbrmstr/ggalt}\AsIs{}
\item[BugReports]\AsIs{}\url{https://github.com/hrbrmstr/ggalt/issues}\AsIs{}
\item[Encoding]\AsIs{UTF-8}
\item[Depends]\AsIs{R (>= 3.5.0), ggplot2 (>= 3.3.0)}
\item[Suggests]\AsIs{testthat, gridExtra, knitr, rmarkdown, ggthemes, reshape2,
maps}
\item[Imports]\AsIs{utils, graphics, datasets, grDevices, dplyr, RColorBrewer,
KernSmooth, proj4, scales, grid, gtable, ash, MASS, extrafont,
tibble, plotly (>= 3.4.1)}
\item[RoxygenNote]\AsIs{7.2.3}
\item[VignetteBuilder]\AsIs{knitr}
\item[Collate]\AsIs{'annotate_textp.r' 'annotation_ticks.r' 'coord_proj.r'
'formatters.r' 'fortify.r' 'position-dodgev.R' 'geom2plotly.r'
'geom_ash.r' 'geom_bkde.r' 'geom_bkde2d.r' 'geom_spikelines.R'
'geom_dumbbell.R' 'geom_cartogram.r' 'geom_encircle.r'
'geom_horizon.r' 'geom_lollipop.r' 'geom_table.r'
'geom_twoway_bar.r' 'geom_xspline.r' 'geom_xspline2.r'
'geom_ubar.r' 'stat-stepribbon.r' 'ggalt-package.r'
'grob_absolute.r' 'guide_axis.r' 'stateface.r' 'utils.r'
'zzz.r'}
\item[NeedsCompilation]\AsIs{no}
\item[Author]\AsIs{Bob Rudis [aut, cre] (ORCID: <}\url{https://orcid.org/0000-0001-5670-2640}\AsIs{>),
Ben Bolker [aut, ctb] (Encircling & additional splines),
Ben Marwick [ctb] (General codebase cleanup),
Jan Schulz [aut, ctb] (Annotations),
Rosen Matev [ctb] (Original annotate_textp implementation on
  stackoverflow),
ProPublica [dtc] (StateFace font),
Aditya Kothari [aut, ctb] (Core functionality of horizon plots),
Ather [dtc] (Core functionality of horizon plots),
Jonathan Sidi [aut, ctb] (Annotation ticks),
Tarcisio Fedrizzi [ctb] (Bytes formatter)}
\end{description}
\Rdcontents{Contents}
\makeatletter\ifthenelse{\boolean{Rd@graphicspath@needs@quotes}}{\graphicspath{{"/Users/mwaite3/Documents/DataExperiments/ggalt/ggalt.Rcheck/ggalt/help/figures/"}}}{\graphicspath{{/Users/mwaite3/Documents/DataExperiments/ggalt/ggalt.Rcheck/ggalt/help/figures/}}}\makeatother
\HeaderA{absoluteGrob}{Absolute grob}{absoluteGrob}
\keyword{internal}{absoluteGrob}
%
\begin{Description}
This grob has fixed dimensions and position.
\end{Description}
%
\begin{Usage}
\begin{verbatim}
absoluteGrob(
  grob,
  width = NULL,
  height = NULL,
  xmin = NULL,
  ymin = NULL,
  vp = NULL
)
\end{verbatim}
\end{Usage}
%
\begin{Details}
It's still experimental
\end{Details}
\HeaderA{annotate\_textp}{Text annotations in plot coordinate system}{annotate.Rul.textp}
%
\begin{Description}
Annotates the plot with text. Compared to \code{annotate("text",...)}, the
placement of the annotations is specified in plot coordinates (from 0 to 1)
instead of data coordinates.
\end{Description}
%
\begin{Usage}
\begin{verbatim}
annotate_textp(
  label,
  x,
  y,
  facets = NULL,
  hjust = 0,
  vjust = 0,
  color = "black",
  alpha = NA,
  family = theme_get()$text$family,
  size = theme_get()$text$size,
  fontface = 1,
  lineheight = 1,
  box_just = ifelse(c(x, y) < 0.5, 0, 1),
  margin = unit(size/2, "pt")
)
\end{verbatim}
\end{Usage}
%
\begin{Arguments}
\begin{ldescription}
\item[\code{label}] text annotation to be placed on the plot

\item[\code{x}, \code{y}] positions of the individual annotations, in plot coordinates
(0..1) instead of data coordinates!

\item[\code{facets}] facet positions of the individual annotations

\item[\code{hjust}, \code{vjust}] horizontal and vertical justification of the text relative
to the bounding box

\item[\code{color}] alpha, family, size, fontface, lineheight font properties

\item[\code{alpha}, \code{family}, \code{size}, \code{fontface}, \code{lineheight}] standard aesthetic customizations

\item[\code{box\_just}] placement of the bounding box for the text relative to x,y
coordinates. Per default, the box is placed to the center of the plot. Be
aware that parts of the box which are outside of the visible region of the
plot will not be shown.

\item[\code{margin}] margins of the bounding box
\end{ldescription}
\end{Arguments}
%
\begin{Examples}
\begin{ExampleCode}
p <- ggplot(mtcars, aes(x = wt, y = mpg)) + geom_point()
p <- p + geom_smooth(method = "lm", se = FALSE)
p + annotate_textp(x = 0.9, y = 0.35, label="A relative linear\nrelationship", hjust=1, color="red")
\end{ExampleCode}
\end{Examples}
\HeaderA{annotation\_ticks}{Annotation: tick marks}{annotation.Rul.ticks}
%
\begin{Description}
This annotation adds tick marks to an axis
\end{Description}
%
\begin{Usage}
\begin{verbatim}
annotation_ticks(
  sides = "b",
  scale = "identity",
  scaled = TRUE,
  short = unit(0.1, "cm"),
  mid = unit(0.2, "cm"),
  long = unit(0.3, "cm"),
  colour = "black",
  size = 0.5,
  linetype = 1,
  alpha = 1,
  color = NULL,
  ticks_per_base = NULL,
  ...
)
\end{verbatim}
\end{Usage}
%
\begin{Arguments}
\begin{ldescription}
\item[\code{sides}] a string that controls which sides of the plot the log ticks appear on.
It can be set to a string containing any of \code{"trbl"}, for top, right,
bottom, and left.

\item[\code{scale}] character, vector of type of scale attributed to each corresponding side, Default: 'identity'

\item[\code{scaled}] is the data already log-scaled? This should be \code{TRUE}
(default) when the data is already transformed with \code{log10()} or when
using \code{scale\_y\_log10()}. It should be \code{FALSE} when using
\code{coord\_trans(y = "log10")}.

\item[\code{short}] a \code{\LinkA{grid::unit()}{grid::unit()}} object specifying the length of the
short tick marks

\item[\code{mid}] a \code{\LinkA{grid::unit()}{grid::unit()}} object specifying the length of the
middle tick marks. In base 10, these are the "5" ticks.

\item[\code{long}] a \code{\LinkA{grid::unit()}{grid::unit()}} object specifying the length of the
long tick marks. In base 10, these are the "1" (or "10") ticks.

\item[\code{colour}] Colour of the tick marks.

\item[\code{size}] Thickness of tick marks, in mm.

\item[\code{linetype}] Linetype of tick marks (\code{solid}, \code{dashed}, etc.)

\item[\code{alpha}] The transparency of the tick marks.

\item[\code{color}] An alias for \code{colour}.

\item[\code{ticks\_per\_base}] integer, number of minor ticks between each pair of major ticks, Default: NULL

\item[\code{...}] Other parameters passed on to the layer
\end{ldescription}
\end{Arguments}
%
\begin{Details}
If scale is of length one it will be replicated to the number of sides given, but if the
length of scale is larger than one it must match the number of sides given.
If ticks\_per\_base is set to NULL the function infers the number of ticks per base to be
the base of the scale - 1, for example log scale is base exp(1) and
log10 and identity are base 10. If ticks\_per\_base is given it follows the same logic as scale.
\end{Details}
%
\begin{Author}
Jonathan Sidi
\end{Author}
%
\begin{Examples}
\begin{ExampleCode}

p <- ggplot(msleep, aes(bodywt, brainwt)) + geom_point()

# Default behavior

# add identity scale minor ticks on y axis
p + annotation_ticks(sides = 'l')

# add identity scale minor ticks on x,y axis
p + annotation_ticks(sides = 'lb')

# Control number of minor ticks of each side independently

# add identity scale minor ticks on x,y axis
p + annotation_ticks(sides = 'lb', ticks_per_base = c(10,5))

# log10 scale
p1 <- p + scale_x_log10()

# add minor ticks on log10 scale
p1 + annotation_ticks(sides = 'b', scale = 'log10')

# add minor ticks on both scales
p1 + annotation_ticks(sides = 'lb', scale = c('identity','log10'))

# add minor ticks on both scales, but force x axis to be identity
p1 + annotation_ticks(sides = 'lb', scale = 'identity')

# log scale
p2 <- p + scale_x_continuous(trans = 'log')

# add minor ticks on log scale
p2 + annotation_ticks(sides = 'b', scale = 'log')

# add minor ticks on both scales
p2 + annotation_ticks(sides = 'lb', scale = c('identity','log'))

# add minor ticks on both scales, but force x axis to be identity
p2 + annotation_ticks(sides = 'lb', scale = 'identity')

\end{ExampleCode}
\end{Examples}
\HeaderA{byte\_format}{Bytes formatter: convert to byte measurement and display symbol.}{byte.Rul.format}
\aliasA{bytes}{byte\_format}{bytes}
\aliasA{Gb}{byte\_format}{Gb}
\aliasA{Kb}{byte\_format}{Kb}
\aliasA{Mb}{byte\_format}{Mb}
%
\begin{Description}
Bytes formatter: convert to byte measurement and display symbol.
\end{Description}
%
\begin{Usage}
\begin{verbatim}
byte_format(symbol = "auto", units = "binary", only_highest = TRUE)

Kb(x)

Mb(x)

Gb(x)

bytes(x, symbol = "auto", units = c("binary", "si"), only_highest = FALSE)
\end{verbatim}
\end{Usage}
%
\begin{Arguments}
\begin{ldescription}
\item[\code{symbol}] byte symbol to use. If "\code{auto}" the symbol used will be
determined by the maximum value of \code{x}. Valid symbols are
"\code{b}", "\code{K}", "\code{Mb}", "\code{Gb}", "\code{Tb}", "\code{Pb}",
"\code{Eb}", "\code{Zb}", and "\code{Yb}", along with their upper case
equivalents and "\code{iB}" equivalents.

\item[\code{units}] which unit base to use, "\code{binary}" (1024 base) or
"\code{si}" (1000 base) for ISI units.

\item[\code{only\_highest}] Whether to use the unit of the highest number or
each number uses its own unit.

\item[\code{x}] a numeric vector to format
\end{ldescription}
\end{Arguments}
%
\begin{Value}
a function with three parameters, \code{x}, a numeric vector that
returns a character vector, \code{symbol} a single or a vector of byte
symbol(s) (e.g. "\code{Kb}") desired and the measurement \code{units}
(traditional \code{binary} or \code{si} for ISI metric units).
\end{Value}
%
\begin{References}
Units of Information (Wikipedia) :
\url{http://en.wikipedia.org/wiki/Units_of_information}
\end{References}
%
\begin{Examples}
\begin{ExampleCode}
byte_format()(sample(3000000000, 10))
bytes(sample(3000000000, 10))
Kb(sample(3000000000, 10))
Mb(sample(3000000000, 10))
Gb(sample(3000000000, 10))
\end{ExampleCode}
\end{Examples}
\HeaderA{CoordProj}{Geom Proto}{CoordProj}
\aliasA{GeomBkde}{CoordProj}{GeomBkde}
\aliasA{GeomBkde2d}{CoordProj}{GeomBkde2d}
\aliasA{GeomDumbbell}{CoordProj}{GeomDumbbell}
\aliasA{GeomEncircle}{CoordProj}{GeomEncircle}
\aliasA{GeomLollipop}{CoordProj}{GeomLollipop}
\aliasA{GeomSpikelines}{CoordProj}{GeomSpikelines}
\aliasA{GeomStateface}{CoordProj}{GeomStateface}
\aliasA{GeomUbar}{CoordProj}{GeomUbar}
\aliasA{GeomXspline}{CoordProj}{GeomXspline}
\aliasA{GeomXSpline2}{CoordProj}{GeomXSpline2}
\aliasA{StatAsh}{CoordProj}{StatAsh}
\aliasA{StatBkde}{CoordProj}{StatBkde}
\aliasA{StatBkde2d}{CoordProj}{StatBkde2d}
\aliasA{StatStepribbon}{CoordProj}{StatStepribbon}
\aliasA{StatXspline}{CoordProj}{StatXspline}
\keyword{datasets}{CoordProj}
\keyword{internal}{CoordProj}
%
\begin{Description}
Geom Proto

Geom Proto

Geom Proto

Geom Proto

Geom Proto
\end{Description}
%
\begin{References}
\url{https://groups.google.com/forum/?fromgroups=\#!topic/ggplot2/9cFWHaH1CPs}
\end{References}
\HeaderA{coord\_proj}{Similar to \code{coord\_map} but uses the PROJ.4 library/package for projection transformation}{coord.Rul.proj}
%
\begin{Description}
The representation of a portion of the earth, which is approximately
spherical, onto a flat 2D plane requires a projection. This is what
\code{coord\_proj} does, using the \code{proj4::project()} function from
the \code{proj4} package.
\end{Description}
%
\begin{Usage}
\begin{verbatim}
coord_proj(
  proj = NULL,
  inverse = FALSE,
  degrees = TRUE,
  ellps.default = "sphere",
  xlim = NULL,
  ylim = NULL
)
\end{verbatim}
\end{Usage}
%
\begin{Arguments}
\begin{ldescription}
\item[\code{proj}] projection definition. If left \code{NULL} will default to
a Robinson projection

\item[\code{inverse}] if \code{TRUE} inverse projection is performed (from a
cartographic projection into lat/long), otherwise projects from
lat/long into a cartographic projection.

\item[\code{degrees}] if \code{TRUE} then the lat/long data is assumed to be in
degrees, otherwise in radians

\item[\code{ellps.default}] default ellipsoid that will be added if no datum or
ellipsoid parameter is specified in proj. Older versions of PROJ.4
didn't require a datum (and used sphere by default), but 4.5.0 and
higher always require a datum or an ellipsoid. Set to  \code{NA} if no
datum should be added to proj (e.g. if you specify an ellipsoid
directly).

\item[\code{xlim}] manually specify x limits (in degrees of longitude)

\item[\code{ylim}] manually specify y limits (in degrees of latitude)
\end{ldescription}
\end{Arguments}
%
\begin{Details}



A sample of the output from \code{coord\_proj()} using the Winkel-Tripel projection:
``
\Figure{coordproj01.png}{width=10cm}

\end{Details}
%
\begin{Note}
It is recommended that you use \code{geom\_cartogram} with this coordinate system

When \code{inverse} is \code{FALSE} \code{coord\_proj} makes a fairly
large assumption that the coordinates being transformed are within
-180:180 (longitude) and -90:90 (latitude). As such, it truncates
all longitude \& latitude input to fit within these ranges. More updates
to this new \code{coord\_} are planned.
\end{Note}
%
\begin{Examples}
\begin{ExampleCode}
## Not run: 
# World in Winkel-Tripel

# U.S.A. Albers-style
usa <- world[world$region == "USA",]
usa <- usa[!(usa$subregion %in% c("Alaska", "Hawaii")),]

gg <- ggplot()
gg <- gg + geom_cartogram(data=usa, map=usa,
                    aes(x=long, y=lat, map_id=region))
gg <- gg + coord_proj(
             paste0("+proj=aea +lat_1=29.5 +lat_2=45.5 +lat_0=37.5 +lon_0=-96",
                    " +x_0=0 +y_0=0 +ellps=GRS80 +datum=NAD83 +units=m +no_defs"))
gg

# Showcase Greenland (properly)
greenland <- world[world$region == "Greenland",]

gg <- ggplot()
gg <- gg + geom_cartogram(data=greenland, map=greenland,
                    aes(x=long, y=lat, map_id=region))
gg <- gg + coord_proj(
             paste0("+proj=stere +lat_0=90 +lat_ts=70 +lon_0=-45 +k=1 +x_0=0",
                    " +y_0=0 +ellps=WGS84 +datum=WGS84 +units=m +no_defs"))
gg

## End(Not run)
\end{ExampleCode}
\end{Examples}
\HeaderA{fortify.table}{Fortify contingency tables}{fortify.table}
%
\begin{Description}
Fortify contingency tables
\end{Description}
%
\begin{Usage}
\begin{verbatim}
## S3 method for class 'table'
fortify(model, data, ...)
\end{verbatim}
\end{Usage}
%
\begin{Arguments}
\begin{ldescription}
\item[\code{model}] the contingency table

\item[\code{data}] data (unused)

\item[\code{...}] (unused)
\end{ldescription}
\end{Arguments}
\HeaderA{GeomTicks}{Base ggproto classes for ggplot2}{GeomTicks}
\aliasA{GeomCartogram}{GeomTicks}{GeomCartogram}
\keyword{datasets}{GeomTicks}
%
\begin{Description}
If you are creating a new geom, stat, position, or scale in another package,
you'll need to extend from ggplot2::Geom, ggplot2::Stat, ggplot2::Position, or ggplot2::Scale.
\end{Description}
%
\begin{SeeAlso}
\code{\LinkA{ggplot2-ggproto}{ggplot2.Rdash.ggproto}}
\end{SeeAlso}
\HeaderA{geom\_bkde}{Display a smooth density estimate.}{geom.Rul.bkde}
\aliasA{stat\_bkde}{geom\_bkde}{stat.Rul.bkde}
%
\begin{Description}
A kernel density estimate, useful for displaying the distribution of
variables with underlying smoothness.
\end{Description}
%
\begin{Usage}
\begin{verbatim}
geom_bkde(
  mapping = NULL,
  data = NULL,
  stat = "bkde",
  position = "identity",
  bandwidth = NULL,
  range.x = NULL,
  na.rm = FALSE,
  show.legend = NA,
  inherit.aes = TRUE,
  ...
)

stat_bkde(
  mapping = NULL,
  data = NULL,
  geom = "area",
  position = "stack",
  kernel = "normal",
  canonical = FALSE,
  bandwidth = NULL,
  gridsize = 410,
  range.x = NULL,
  truncate = TRUE,
  na.rm = FALSE,
  show.legend = NA,
  inherit.aes = TRUE,
  ...
)
\end{verbatim}
\end{Usage}
%
\begin{Arguments}
\begin{ldescription}
\item[\code{mapping}] Set of aesthetic mappings created by \code{\LinkA{aes()}{aes()}}. If specified and
\code{inherit.aes = TRUE} (the default), it is combined with the default mapping
at the top level of the plot. You must supply \code{mapping} if there is no plot
mapping.

\item[\code{data}] The data to be displayed in this layer. There are three
options:

If \code{NULL}, the default, the data is inherited from the plot
data as specified in the call to \code{\LinkA{ggplot()}{ggplot()}}.

A \code{data.frame}, or other object, will override the plot
data. All objects will be fortified to produce a data frame. See
\code{\LinkA{fortify()}{fortify()}} for which variables will be created.

A \code{function} will be called with a single argument,
the plot data. The return value must be a \code{data.frame}, and
will be used as the layer data. A \code{function} can be created
from a \code{formula} (e.g. \code{\textasciitilde{} head(.x, 10)}).

\item[\code{position}] Position adjustment, either as a string naming the adjustment
(e.g. \code{"jitter"} to use \code{position\_jitter}), or the result of a call to a
position adjustment function. Use the latter if you need to change the
settings of the adjustment.

\item[\code{bandwidth}] the kernel bandwidth smoothing parameter. see
\code{\LinkA{bkde}{bkde}} for details. If \code{NULL},
it will be computed for you but will most likely not yield optimal
results.

\item[\code{range.x}] vector containing the minimum and maximum values of x at which
to compute the estimate. see \code{\LinkA{bkde}{bkde}} for details

\item[\code{na.rm}] If \code{FALSE}, the default, missing values are removed with
a warning. If \code{TRUE}, missing values are silently removed.

\item[\code{show.legend}] logical. Should this layer be included in the legends?
\code{NA}, the default, includes if any aesthetics are mapped.
\code{FALSE} never includes, and \code{TRUE} always includes.
It can also be a named logical vector to finely select the aesthetics to
display.

\item[\code{inherit.aes}] If \code{FALSE}, overrides the default aesthetics,
rather than combining with them. This is most useful for helper functions
that define both data and aesthetics and shouldn't inherit behaviour from
the default plot specification, e.g. \code{\LinkA{borders()}{borders()}}.

\item[\code{...}] Other arguments passed on to \code{\LinkA{layer()}{layer()}}. These are
often aesthetics, used to set an aesthetic to a fixed value, like
\code{colour = "red"} or \code{size = 3}. They may also be parameters
to the paired geom/stat.

\item[\code{geom}, \code{stat}] Use to override the default connection between
\code{geom\_bkde} and \code{stat\_bkde}.

\item[\code{kernel}] character string which determines the smoothing kernel. see
\code{\LinkA{bkde}{bkde}} for details

\item[\code{canonical}] logical flag: if TRUE, canonically scaled kernels are used.
see \code{\LinkA{bkde}{bkde}} for details

\item[\code{gridsize}] the number of equally spaced points at which to estimate the
density. see \code{\LinkA{bkde}{bkde}} for details.

\item[\code{truncate}] logical flag: if TRUE, data with x values outside the range
specified by range.x are ignored. see \code{\LinkA{bkde}{bkde}}
for details
\end{ldescription}
\end{Arguments}
%
\begin{Details}



A sample of the output from \code{geom\_bkde()}:

\Figure{geombkde01.png}{width=10cm}

\end{Details}
%
\begin{Section}{Aesthetics}

\code{geom\_bkde} understands the following aesthetics (required aesthetics
are in bold):
\begin{itemize}

\item{} \strong{\code{x}}
\item{} \strong{\code{y}}
\item{} \code{alpha}
\item{} \code{color}
\item{} \code{fill}
\item{} \code{linetype}
\item{} \code{size}

\end{itemize}

\end{Section}
%
\begin{Section}{Computed variables}

\begin{description}

\item[density] density estimate
\item[count] density * number of points - useful for stacked density
plots
\item[scaled] density estimate, scaled to maximum of 1

\end{description}

\end{Section}
%
\begin{SeeAlso}
See \code{\LinkA{geom\_histogram}{geom.Rul.histogram}}, \code{\LinkA{geom\_freqpoly}{geom.Rul.freqpoly}} for
other methods of displaying continuous distribution.
See \code{\LinkA{geom\_violin}{geom.Rul.violin}} for a compact density display.
\end{SeeAlso}
%
\begin{Examples}
\begin{ExampleCode}
data(geyser, package="MASS")

ggplot(geyser, aes(x=duration)) +
  stat_bkde(alpha=1/2)

ggplot(geyser, aes(x=duration)) +
  geom_bkde(alpha=1/2)

ggplot(geyser, aes(x=duration)) +
 stat_bkde(bandwidth=0.25)

ggplot(geyser, aes(x=duration)) +
  geom_bkde(bandwidth=0.25)
\end{ExampleCode}
\end{Examples}
\HeaderA{geom\_bkde2d}{Contours from a 2d density estimate.}{geom.Rul.bkde2d}
\aliasA{stat\_bkde2d}{geom\_bkde2d}{stat.Rul.bkde2d}
%
\begin{Description}
Perform a 2D kernel density estimation using \code{bkde2D} and display the
results with contours. This can be useful for dealing with overplotting
\end{Description}
%
\begin{Usage}
\begin{verbatim}
geom_bkde2d(
  mapping = NULL,
  data = NULL,
  stat = "bkde2d",
  position = "identity",
  bandwidth = NULL,
  range.x = NULL,
  lineend = "butt",
  contour = TRUE,
  linejoin = "round",
  linemitre = 1,
  na.rm = FALSE,
  show.legend = NA,
  inherit.aes = TRUE,
  ...
)

stat_bkde2d(
  mapping = NULL,
  data = NULL,
  geom = "density2d",
  position = "identity",
  contour = TRUE,
  bandwidth = NULL,
  grid_size = c(51, 51),
  range.x = NULL,
  truncate = TRUE,
  na.rm = FALSE,
  show.legend = NA,
  inherit.aes = TRUE,
  ...
)
\end{verbatim}
\end{Usage}
%
\begin{Arguments}
\begin{ldescription}
\item[\code{mapping}] Set of aesthetic mappings created by \code{\LinkA{aes()}{aes()}}. If specified and
\code{inherit.aes = TRUE} (the default), it is combined with the default mapping
at the top level of the plot. You must supply \code{mapping} if there is no plot
mapping.

\item[\code{data}] The data to be displayed in this layer. There are three
options:

If \code{NULL}, the default, the data is inherited from the plot
data as specified in the call to \code{\LinkA{ggplot()}{ggplot()}}.

A \code{data.frame}, or other object, will override the plot
data. All objects will be fortified to produce a data frame. See
\code{\LinkA{fortify()}{fortify()}} for which variables will be created.

A \code{function} will be called with a single argument,
the plot data. The return value must be a \code{data.frame}, and
will be used as the layer data. A \code{function} can be created
from a \code{formula} (e.g. \code{\textasciitilde{} head(.x, 10)}).

\item[\code{stat}] The statistical transformation to use on the data for this
layer, either as a \code{ggproto} \code{Geom} subclass or as a string naming the
stat stripped of the \code{stat\_} prefix (e.g. \code{"count"} rather than
\code{"stat\_count"})

\item[\code{position}] Position adjustment, either as a string naming the adjustment
(e.g. \code{"jitter"} to use \code{position\_jitter}), or the result of a call to a
position adjustment function. Use the latter if you need to change the
settings of the adjustment.

\item[\code{bandwidth}] the kernel bandwidth smoothing parameter. see
\code{\LinkA{bkde2D}{bkde2D}} for details. If \code{NULL},
it will be computed for you but will most likely not yield optimal
results. see \code{\LinkA{bkde2D}{bkde2D}} for details

\item[\code{range.x}] a list containing two vectors, where each vector contains the
minimum and maximum values of x at which to compute the estimate for
each direction. see \code{\LinkA{bkde2D}{bkde2D}} for details

\item[\code{lineend}] Line end style (round, butt, square).

\item[\code{contour}] If \code{TRUE}, contour the results of the 2d density
estimation

\item[\code{linejoin}] Line join style (round, mitre, bevel).

\item[\code{linemitre}] Line mitre limit (number greater than 1).

\item[\code{na.rm}] If \code{FALSE}, the default, missing values are removed with
a warning. If \code{TRUE}, missing values are silently removed.

\item[\code{show.legend}] logical. Should this layer be included in the legends?
\code{NA}, the default, includes if any aesthetics are mapped.
\code{FALSE} never includes, and \code{TRUE} always includes.
It can also be a named logical vector to finely select the aesthetics to
display.

\item[\code{inherit.aes}] If \code{FALSE}, overrides the default aesthetics,
rather than combining with them. This is most useful for helper functions
that define both data and aesthetics and shouldn't inherit behaviour from
the default plot specification, e.g. \code{\LinkA{borders()}{borders()}}.

\item[\code{...}] Other arguments passed on to \code{\LinkA{layer()}{layer()}}. These are
often aesthetics, used to set an aesthetic to a fixed value, like
\code{colour = "red"} or \code{size = 3}. They may also be parameters
to the paired geom/stat.

\item[\code{geom}] default geom to use with this stat

\item[\code{grid\_size}] vector containing the number of equally spaced points in each
direction over which the density is to be estimated. see
\code{\LinkA{bkde2D}{bkde2D}} for details

\item[\code{truncate}] logical flag: if TRUE, data with x values outside the range
specified by range.x are ignored. see \code{\LinkA{bkde2D}{bkde2D}}
for details
\end{ldescription}
\end{Arguments}
%
\begin{Details}



A sample of the output from \code{geom\_bkde2d()}:

\Figure{geombkde2d01.png}{width=10cm}

\end{Details}
%
\begin{Section}{Computed variables}

Same as \code{\LinkA{stat\_contour}{stat.Rul.contour}}
\end{Section}
%
\begin{SeeAlso}
\code{\LinkA{geom\_contour}{geom.Rul.contour}} for contour drawing geom,
\code{\LinkA{stat\_sum}{stat.Rul.sum}} for another way of dealing with overplotting
\end{SeeAlso}
%
\begin{Examples}
\begin{ExampleCode}
m <- ggplot(faithful, aes(x = eruptions, y = waiting)) +
       geom_point() +
       xlim(0.5, 6) +
       ylim(40, 110)

m + geom_bkde2d(bandwidth=c(0.5, 4))

m + stat_bkde2d(bandwidth=c(0.5, 4), aes(fill = ..level..), geom = "polygon")

# If you map an aesthetic to a categorical variable, you will get a
# set of contours for each value of that variable
set.seed(4393)
dsmall <- diamonds[sample(nrow(diamonds), 1000), ]
d <- ggplot(dsmall, aes(x, y)) +
       geom_bkde2d(bandwidth=c(0.5, 0.5), aes(colour = cut))
d

# If we turn contouring off, we can use use geoms like tiles:
d + stat_bkde2d(bandwidth=c(0.5, 0.5), geom = "raster",
                aes(fill = ..density..), contour = FALSE)

# Or points:
d + stat_bkde2d(bandwidth=c(0.5, 0.5), geom = "point",
                aes(size = ..density..),  contour = FALSE)
\end{ExampleCode}
\end{Examples}
\HeaderA{geom\_cartogram}{Map polygons layer enabling the display of show statistical information}{geom.Rul.cartogram}
%
\begin{Description}
This replicates the old behaviour of \code{geom\_map()}, enabling specifying of
\code{x} and \code{y} aesthetics.
\end{Description}
%
\begin{Usage}
\begin{verbatim}
geom_cartogram(
  mapping = NULL,
  data = NULL,
  stat = "identity",
  ...,
  map,
  na.rm = FALSE,
  show.legend = NA,
  inherit.aes = TRUE
)
\end{verbatim}
\end{Usage}
%
\begin{Arguments}
\begin{ldescription}
\item[\code{mapping}] Set of aesthetic mappings created by \code{\LinkA{aes()}{aes()}}. If specified and
\code{inherit.aes = TRUE} (the default), it is combined with the default mapping
at the top level of the plot. You must supply \code{mapping} if there is no plot
mapping.

\item[\code{data}] The data to be displayed in this layer. There are three
options:

If \code{NULL}, the default, the data is inherited from the plot
data as specified in the call to \code{\LinkA{ggplot()}{ggplot()}}.

A \code{data.frame}, or other object, will override the plot
data. All objects will be fortified to produce a data frame. See
\code{\LinkA{fortify()}{fortify()}} for which variables will be created.

A \code{function} will be called with a single argument,
the plot data. The return value must be a \code{data.frame}, and
will be used as the layer data. A \code{function} can be created
from a \code{formula} (e.g. \code{\textasciitilde{} head(.x, 10)}).

\item[\code{stat}] The statistical transformation to use on the data for this
layer, either as a \code{ggproto} \code{Geom} subclass or as a string naming the
stat stripped of the \code{stat\_} prefix (e.g. \code{"count"} rather than
\code{"stat\_count"})

\item[\code{...}] Other arguments passed on to \code{\LinkA{layer()}{layer()}}. These are
often aesthetics, used to set an aesthetic to a fixed value, like
\code{colour = "red"} or \code{size = 3}. They may also be parameters
to the paired geom/stat.

\item[\code{map}] Data frame that contains the map coordinates.  This will
typically be created using \code{\LinkA{fortify}{fortify}} on a spatial object.
It must contain columns \code{x}, \code{long} or \code{longitude},
\code{y}, \code{lat} or \code{latitude} and \code{region} or \code{id}.

\item[\code{na.rm}] If \code{FALSE}, the default, missing values are removed with
a warning. If \code{TRUE}, missing values are silently removed.

\item[\code{show.legend}] logical. Should this layer be included in the legends?
\code{NA}, the default, includes if any aesthetics are mapped.
\code{FALSE} never includes, and \code{TRUE} always includes.
It can also be a named logical vector to finely select the aesthetics to
display.

\item[\code{inherit.aes}] If \code{FALSE}, overrides the default aesthetics,
rather than combining with them. This is most useful for helper functions
that define both data and aesthetics and shouldn't inherit behaviour from
the default plot specification, e.g. \code{\LinkA{borders()}{borders()}}.
\end{ldescription}
\end{Arguments}
%
\begin{Section}{Aesthetics}

\code{geom\_cartogram} understands the following aesthetics (required aesthetics are in bold):
\begin{itemize}

\item{} \code{map\_id}
\item{} \code{alpha}
\item{} \code{colour}
\item{} \code{fill}
\item{} \code{group}
\item{} \code{linetype}
\item{} \code{size}
\item{} \code{x}
\item{} \code{y}

\end{itemize}

\end{Section}
%
\begin{Examples}
\begin{ExampleCode}
## Not run: 
# When using geom_polygon, you will typically need two data frames:
# one contains the coordinates of each polygon (positions),  and the
# other the values associated with each polygon (values).  An id
# variable links the two together

ids <- factor(c("1.1", "2.1", "1.2", "2.2", "1.3", "2.3"))

values <- data.frame(
  id = ids,
  value = c(3, 3.1, 3.1, 3.2, 3.15, 3.5)
)

positions <- data.frame(
  id = rep(ids, each = 4),
  x = c(2, 1, 1.1, 2.2, 1, 0, 0.3, 1.1, 2.2, 1.1, 1.2, 2.5, 1.1, 0.3,
  0.5, 1.2, 2.5, 1.2, 1.3, 2.7, 1.2, 0.5, 0.6, 1.3),
  y = c(-0.5, 0, 1, 0.5, 0, 0.5, 1.5, 1, 0.5, 1, 2.1, 1.7, 1, 1.5,
  2.2, 2.1, 1.7, 2.1, 3.2, 2.8, 2.1, 2.2, 3.3, 3.2)
)

ggplot() +
  geom_cartogram(aes(x, y, map_id = id), map = positions, data=positions)

ggplot() +
  geom_cartogram(aes(x, y, map_id = id), map = positions, data=positions) +
  geom_cartogram(data=values, map=positions, aes(fill = value, map_id=id))

ggplot() +
  geom_cartogram(aes(x, y, map_id = id), map = positions, data=positions) +
  geom_cartogram(data=values, map=positions, aes(fill = value, map_id=id)) +
  ylim(0, 3)

# Better example
crimes <- data.frame(state = tolower(rownames(USArrests)), USArrests)
crimesm <- reshape2::melt(crimes, id = 1)

if (require(maps)) {

  states_map <- map_data("state")

  ggplot() +
    geom_cartogram(aes(long, lat, map_id = region), map = states_map, data=states_map) +
    geom_cartogram(aes(fill = Murder, map_id = state), map=states_map, data=crimes)

  last_plot() + coord_map("polyconic")

  ggplot() +
    geom_cartogram(aes(long, lat, map_id=region), map = states_map, data=states_map) +
    geom_cartogram(aes(fill = value, map_id=state), map = states_map, data=crimesm) +
    coord_map("polyconic") +
    facet_wrap( ~ variable)
}

## End(Not run)
\end{ExampleCode}
\end{Examples}
\HeaderA{geom\_dumbbell}{Dumbbell charts}{geom.Rul.dumbbell}
%
\begin{Description}
The dumbbell geom is used to create dumbbell charts.
\end{Description}
%
\begin{Usage}
\begin{verbatim}
geom_dumbbell(
  mapping = NULL,
  data = NULL,
  ...,
  colour_x = NULL,
  size_x = NULL,
  colour_xend = NULL,
  size_xend = NULL,
  dot_guide = FALSE,
  dot_guide_size = NULL,
  dot_guide_colour = NULL,
  na.rm = FALSE,
  show.legend = NA,
  inherit.aes = TRUE,
  position = "identity"
)
\end{verbatim}
\end{Usage}
%
\begin{Arguments}
\begin{ldescription}
\item[\code{mapping}] Set of aesthetic mappings created by \code{\LinkA{aes()}{aes()}}. If specified and
\code{inherit.aes = TRUE} (the default), it is combined with the default mapping
at the top level of the plot. You must supply \code{mapping} if there is no plot
mapping.

\item[\code{data}] The data to be displayed in this layer. There are three
options:

If \code{NULL}, the default, the data is inherited from the plot
data as specified in the call to \code{\LinkA{ggplot()}{ggplot()}}.

A \code{data.frame}, or other object, will override the plot
data. All objects will be fortified to produce a data frame. See
\code{\LinkA{fortify()}{fortify()}} for which variables will be created.

A \code{function} will be called with a single argument,
the plot data. The return value must be a \code{data.frame}, and
will be used as the layer data. A \code{function} can be created
from a \code{formula} (e.g. \code{\textasciitilde{} head(.x, 10)}).

\item[\code{...}] other arguments passed on to \code{\LinkA{layer}{layer}}. These are
often aesthetics, used to set an aesthetic to a fixed value, like
\code{color = "red"} or \code{size = 3}. They may also be parameters
to the paired geom/stat.

\item[\code{colour\_x}] the colour of the start point

\item[\code{size\_x}] the size of the start point

\item[\code{colour\_xend}] the colour of the end point

\item[\code{size\_xend}] the size of the end point

\item[\code{dot\_guide}] if \code{TRUE}, a leading dotted line will be placed before the left-most dumbbell point

\item[\code{dot\_guide\_size}, \code{dot\_guide\_colour}] singe-value aesthetics for \code{dot\_guide}

\item[\code{na.rm}] If \code{FALSE} (the default), removes missing values with
a warning.  If \code{TRUE} silently removes missing values.

\item[\code{show.legend}] logical. Should this layer be included in the legends?
\code{NA}, the default, includes if any aesthetics are mapped.
\code{FALSE} never includes, and \code{TRUE} always includes.
It can also be a named logical vector to finely select the aesthetics to
display.

\item[\code{inherit.aes}] If \code{FALSE}, overrides the default aesthetics,
rather than combining with them. This is most useful for helper functions
that define both data and aesthetics and shouldn't inherit behaviour from
the default plot specification, e.g. \code{\LinkA{borders()}{borders()}}.

\item[\code{position}] Position adjustment, either as a string, or the result of a
call to a position adjustment function.
\end{ldescription}
\end{Arguments}
%
\begin{Details}
Dumbbell dot plots — dot plots with two or more series of data — are an
alternative to the clustered bar chart or slope graph.
\end{Details}
%
\begin{Section}{Aesthetics}

@section Aesthetics:
\code{geom\_segment()}understands the following aesthetics. Required aesthetics are displayed in bold and defaults are displayed for optional aesthetics:

\Tabular{rll}{
• &\strong{\code{\LinkA{x}{x}}}&\\{}
• &\strong{\code{\LinkA{y}{y}}}&\\{}
• &\strong{\code{\LinkA{xend}{xend}}\emph{or}\code{\LinkA{yend}{yend}}}&\\{}
• &\code{\LinkA{alpha}{alpha}}&→ \code{NA}\\{}
• &\code{\LinkA{colour}{colour}}&→ via \code{theme()}\\{}
• &\code{\LinkA{group}{group}}&→ inferred \\{}
• &\code{\LinkA{linetype}{linetype}}&→ via \code{theme()}\\{}
• &\code{\LinkA{linewidth}{linewidth}}&→ via \code{theme()}\\{}
}
Learn more about setting these aesthetics in \code{vignette("ggplot2-specs")}.
\end{Section}
%
\begin{Examples}
\begin{ExampleCode}
library(ggplot2)

df <- data.frame(trt=LETTERS[1:5], l=c(20, 40, 10, 30, 50), r=c(70, 50, 30, 60, 80))

ggplot(df, aes(y=trt, x=l, xend=r)) +
  geom_dumbbell(size=3, color="#e3e2e1",
                colour_x = "#5b8124", colour_xend = "#bad744",
                dot_guide=TRUE, dot_guide_size=0.25) +
  labs(x=NULL, y=NULL, title="ggplot2 geom_dumbbell with dot guide") +
  theme_minimal() +
  theme(panel.grid.major.x=element_line(size=0.05))

## with vertical dodging
df2 <- data.frame(trt = c(LETTERS[1:5], "D"),
                 l = c(20, 40, 10, 30, 50, 40),
                 r = c(70, 50, 30, 60, 80, 70))

ggplot(df2, aes(y=trt, x=l, xend=r)) +
  geom_dumbbell(size=3, color="#e3e2e1",
                colour_x = "#5b8124", colour_xend = "#bad744",
                dot_guide=TRUE, dot_guide_size=0.25,
                position=position_dodgev(height=0.4)) +
  labs(x=NULL, y=NULL, title="ggplot2 geom_dumbbell with dot guide") +
  theme_minimal() +
  theme(panel.grid.major.x=element_line(size=0.05))
\end{ExampleCode}
\end{Examples}
\HeaderA{geom\_encircle}{Automatically enclose points in a polygon}{geom.Rul.encircle}
%
\begin{Description}
Automatically enclose points in a polygon
\end{Description}
%
\begin{Usage}
\begin{verbatim}
geom_encircle(
  mapping = NULL,
  data = NULL,
  stat = "identity",
  position = "identity",
  na.rm = FALSE,
  show.legend = NA,
  inherit.aes = TRUE,
  ...
)
\end{verbatim}
\end{Usage}
%
\begin{Arguments}
\begin{ldescription}
\item[\code{mapping}] mapping

\item[\code{data}] data

\item[\code{stat}] stat

\item[\code{position}] position

\item[\code{na.rm}] na.rm

\item[\code{show.legend}] show.legend

\item[\code{inherit.aes}] inherit.aes

\item[\code{...}] dots
\end{ldescription}
\end{Arguments}
%
\begin{Details}



A sample of the output from \code{geom\_encircle()}:

\Figure{geomencircle01.png}{width=10cm}

\end{Details}
%
\begin{Value}
adds a circle around the specified points
\end{Value}
%
\begin{Author}
Ben Bolker
\end{Author}
%
\begin{Examples}
\begin{ExampleCode}
d <- data.frame(x=c(1,1,2),y=c(1,2,2)*100)

gg <- ggplot(d,aes(x,y))
gg <- gg + scale_x_continuous(expand=c(0.5,1))
gg <- gg + scale_y_continuous(expand=c(0.5,1))

gg + geom_encircle(s_shape=1, expand=0) + geom_point()

gg + geom_encircle(s_shape=1, expand=0.1, colour="red") + geom_point()

gg + geom_encircle(s_shape=0.5, expand=0.1, colour="purple") + geom_point()

gg + geom_encircle(data=subset(d, x==1), colour="blue", spread=0.02) +
  geom_point()

gg +geom_encircle(data=subset(d, x==2), colour="cyan", spread=0.04) +
  geom_point()

gg <- ggplot(mpg, aes(displ, hwy))
gg + geom_encircle(data=subset(mpg, hwy>40)) + geom_point()
gg + geom_encircle(aes(group=manufacturer)) + geom_point()
gg + geom_encircle(aes(group=manufacturer,fill=manufacturer),alpha=0.4)+
       geom_point()
gg + geom_encircle(aes(group=manufacturer,colour=manufacturer))+
       geom_point()

ss <- subset(mpg,hwy>31 & displ<2)

gg + geom_encircle(data=ss, colour="blue", s_shape=0.9, expand=0.07) +
  geom_point() + geom_point(data=ss, colour="blue")
\end{ExampleCode}
\end{Examples}
\HeaderA{geom\_horizon}{Plot a time series as a horizon plot}{geom.Rul.horizon}
\aliasA{GeomHorizon}{geom\_horizon}{GeomHorizon}
\aliasA{StatHorizon}{geom\_horizon}{StatHorizon}
\aliasA{stat\_horizon}{geom\_horizon}{stat.Rul.horizon}
\keyword{internal}{geom\_horizon}
%
\begin{Description}
A horizon plot breaks the Y dimension down using colours. This is useful
when visualising y values spanning a vast range and / or trying to highlight
outliers without losing context of the rest of the data.\\{} \\{} Horizon
plots are best viewed in an apsect ratio of very low vertical length.
\end{Description}
%
\begin{Usage}
\begin{verbatim}
geom_horizon(
  mapping = NULL,
  data = NULL,
  show.legend = TRUE,
  inherit.aes = TRUE,
  na.rm = TRUE,
  bandwidth = NULL,
  ...
)

GeomHorizon

stat_horizon(
  mapping = NULL,
  data = NULL,
  geom = "horizon",
  show.legend = TRUE,
  inherit.aes = TRUE,
  na.rm = TRUE,
  bandwidth = NULL,
  ...
)

StatHorizon
\end{verbatim}
\end{Usage}
%
\begin{Format}
An object of class \code{GeomHorizon} (inherits from \code{GeomArea}, \code{GeomRibbon}, \code{Geom}, \code{ggproto}, \code{gg}) of length 4.

An object of class \code{StatHorizon} (inherits from \code{Stat}, \code{ggproto}, \code{gg}) of length 5.
\end{Format}
%
\begin{Section}{Aesthetics}
\code{x}, \code{y}, \code{fill}. \code{fill} defaults to \code{..band..} which is
the band number the current data fill area belongs in.
\end{Section}
%
\begin{Section}{Other parameters}
\code{bandwidth}, to dictate the span of a band.
\end{Section}
\HeaderA{geom\_lollipop}{Lollipop charts}{geom.Rul.lollipop}
%
\begin{Description}
The lollipop geom is used to create lollipop charts.
\end{Description}
%
\begin{Usage}
\begin{verbatim}
geom_lollipop(
  mapping = NULL,
  data = NULL,
  ...,
  horizontal = FALSE,
  point.colour = NULL,
  point.size = NULL,
  na.rm = FALSE,
  show.legend = NA,
  inherit.aes = TRUE
)
\end{verbatim}
\end{Usage}
%
\begin{Arguments}
\begin{ldescription}
\item[\code{mapping}] Set of aesthetic mappings created by \code{\LinkA{aes()}{aes()}}. If specified and
\code{inherit.aes = TRUE} (the default), it is combined with the default mapping
at the top level of the plot. You must supply \code{mapping} if there is no plot
mapping.

\item[\code{data}] The data to be displayed in this layer. There are three
options:

If \code{NULL}, the default, the data is inherited from the plot
data as specified in the call to \code{\LinkA{ggplot()}{ggplot()}}.

A \code{data.frame}, or other object, will override the plot
data. All objects will be fortified to produce a data frame. See
\code{\LinkA{fortify()}{fortify()}} for which variables will be created.

A \code{function} will be called with a single argument,
the plot data. The return value must be a \code{data.frame}, and
will be used as the layer data. A \code{function} can be created
from a \code{formula} (e.g. \code{\textasciitilde{} head(.x, 10)}).

\item[\code{...}] other arguments passed on to \code{\LinkA{layer}{layer}}. These are
often aesthetics, used to set an aesthetic to a fixed value, like
\code{color = "red"} or \code{size = 3}. They may also be parameters
to the paired geom/stat.

\item[\code{horizontal}] \code{horizontal} is \code{FALSE} (the default), the function
will draw the lollipops up from the X axis (i.e. it will set \code{xend}
to \code{x} \& \code{yend} to \code{0}). If \code{TRUE}, it wiill set
\code{yend} to \code{y} \& \code{xend} to \code{0}). Make sure you map the
\code{x} \& \code{y} aesthetics accordingly. This parameter helps avoid
the need for \code{coord\_flip()}.

\item[\code{point.colour}] the colour of the point

\item[\code{point.size}] the size of the point

\item[\code{na.rm}] If \code{FALSE} (the default), removes missing values with
a warning.  If \code{TRUE} silently removes missing values.

\item[\code{show.legend}] logical. Should this layer be included in the legends?
\code{NA}, the default, includes if any aesthetics are mapped.
\code{FALSE} never includes, and \code{TRUE} always includes.
It can also be a named logical vector to finely select the aesthetics to
display.

\item[\code{inherit.aes}] If \code{FALSE}, overrides the default aesthetics,
rather than combining with them. This is most useful for helper functions
that define both data and aesthetics and shouldn't inherit behaviour from
the default plot specification, e.g. \code{\LinkA{borders()}{borders()}}.
\end{ldescription}
\end{Arguments}
%
\begin{Details}
Lollipop charts are the creation of Andy Cotgreave going back to 2011. They
are a combination of a thin segment, starting at with a dot at the top and are a
suitable alternative to or replacement for bar charts.

Use the \code{horizontal} parameter to abate the need for \code{coord\_flip()}
(see the \code{Arguments} section for details).




A sample of the output from \code{geom\_lollipop()}:

\Figure{geomlollipop01.png}{width=10cm}

\end{Details}
%
\begin{Section}{Aesthetics}

@section Aesthetics:
\code{geom\_point()}understands the following aesthetics. Required aesthetics are displayed in bold and defaults are displayed for optional aesthetics:

\Tabular{rll}{
• &\strong{\code{\LinkA{x}{x}}}&\\{}
• &\strong{\code{\LinkA{y}{y}}}&\\{}
• &\code{\LinkA{alpha}{alpha}}&→ \code{NA}\\{}
• &\code{\LinkA{colour}{colour}}&→ via \code{theme()}\\{}
• &\code{\LinkA{fill}{fill}}&→ via \code{theme()}\\{}
• &\code{\LinkA{group}{group}}&→ inferred \\{}
• &\code{\LinkA{shape}{shape}}&→ via \code{theme()}\\{}
• &\code{\LinkA{size}{size}}&→ via \code{theme()}\\{}
• &\code{stroke}&→ via \code{theme()}\\{}
}
Learn more about setting these aesthetics in \code{vignette("ggplot2-specs")}.
\end{Section}
%
\begin{Examples}
\begin{ExampleCode}
df <- data.frame(trt=LETTERS[1:10],
                 value=seq(100, 10, by=-10))

ggplot(df, aes(trt, value)) + geom_lollipop()

ggplot(df, aes(value, trt)) + geom_lollipop(horizontal=TRUE)
\end{ExampleCode}
\end{Examples}
\HeaderA{geom\_spikelines}{Draw spikelines on a plot}{geom.Rul.spikelines}
%
\begin{Description}
Segment reference lines that originate at an point
\end{Description}
%
\begin{Usage}
\begin{verbatim}
geom_spikelines(
  mapping = NULL,
  data = NULL,
  stat = "identity",
  position = "identity",
  ...,
  arrow = NULL,
  lineend = "butt",
  linejoin = "round",
  na.rm = FALSE,
  show.legend = NA,
  inherit.aes = TRUE
)
\end{verbatim}
\end{Usage}
%
\begin{Arguments}
\begin{ldescription}
\item[\code{mapping}] Set of aesthetic mappings created by \code{\LinkA{aes()}{aes()}}. If specified and
\code{inherit.aes = TRUE} (the default), it is combined with the default mapping
at the top level of the plot. You must supply \code{mapping} if there is no plot
mapping.

\item[\code{data}] The data to be displayed in this layer. There are three
options:

If \code{NULL}, the default, the data is inherited from the plot
data as specified in the call to \code{\LinkA{ggplot()}{ggplot()}}.

A \code{data.frame}, or other object, will override the plot
data. All objects will be fortified to produce a data frame. See
\code{\LinkA{fortify()}{fortify()}} for which variables will be created.

A \code{function} will be called with a single argument,
the plot data. The return value must be a \code{data.frame}, and
will be used as the layer data. A \code{function} can be created
from a \code{formula} (e.g. \code{\textasciitilde{} head(.x, 10)}).

\item[\code{stat}] The statistical transformation to use on the data for this
layer, either as a \code{ggproto} \code{Geom} subclass or as a string naming the
stat stripped of the \code{stat\_} prefix (e.g. \code{"count"} rather than
\code{"stat\_count"})

\item[\code{position}] Position adjustment, either as a string naming the adjustment
(e.g. \code{"jitter"} to use \code{position\_jitter}), or the result of a call to a
position adjustment function. Use the latter if you need to change the
settings of the adjustment.

\item[\code{...}] Other arguments passed on to \code{\LinkA{layer()}{layer()}}. These are
often aesthetics, used to set an aesthetic to a fixed value, like
\code{colour = "red"} or \code{size = 3}. They may also be parameters
to the paired geom/stat.

\item[\code{arrow}] Arrow specification, as created by \code{\LinkA{grid::arrow()}{grid::arrow()}}.

\item[\code{lineend}] Line end style (round, butt, square).

\item[\code{linejoin}] Line join style (round, mitre, bevel).

\item[\code{na.rm}] If \code{FALSE}, the default, missing values are removed with
a warning. If \code{TRUE}, missing values are silently removed.

\item[\code{show.legend}] logical. Should this layer be included in the legends?
\code{NA}, the default, includes if any aesthetics are mapped.
\code{FALSE} never includes, and \code{TRUE} always includes.
It can also be a named logical vector to finely select the aesthetics to
display.

\item[\code{inherit.aes}] If \code{FALSE}, overrides the default aesthetics,
rather than combining with them. This is most useful for helper functions
that define both data and aesthetics and shouldn't inherit behaviour from
the default plot specification, e.g. \code{\LinkA{borders()}{borders()}}.
\end{ldescription}
\end{Arguments}
%
\begin{Author}
Jonathan Sidi
\end{Author}
%
\begin{Examples}
\begin{ExampleCode}

mtcars$name <- rownames(mtcars)

p <- ggplot(data = mtcars, aes(x=mpg,y=disp)) + geom_point()

p + geom_spikelines(data = mtcars[mtcars$carb==4,],linetype = 2)

p + geom_spikelines(data = mtcars[mtcars$carb==4,],aes(colour = factor(gear)), linetype = 2)

## Not run: 
require(ggrepel)
p + geom_spikelines(data = mtcars[mtcars$carb==4,],aes(colour = factor(gear)), linetype = 2) +
ggrepel::geom_label_repel(data = mtcars[mtcars$carb==4,],aes(label = name))

## End(Not run)

\end{ExampleCode}
\end{Examples}
\HeaderA{geom\_stateface}{Use ProPublica's StateFace font in ggplot2 plots}{geom.Rul.stateface}
\keyword{StateFace operations}{geom\_stateface}
%
\begin{Description}
The \code{label} parameter can be either a 2-letter state abbreviation
or a full state name. \code{geom\_stateface()} will take care of the
translation to StateFace font glyph characters.
\end{Description}
%
\begin{Usage}
\begin{verbatim}
geom_stateface(
  mapping = NULL,
  data = NULL,
  stat = "identity",
  position = "identity",
  ...,
  parse = FALSE,
  nudge_x = 0,
  nudge_y = 0,
  check_overlap = FALSE,
  na.rm = FALSE,
  show.legend = NA,
  inherit.aes = TRUE
)
\end{verbatim}
\end{Usage}
%
\begin{Arguments}
\begin{ldescription}
\item[\code{mapping}] Set of aesthetic mappings created by \code{\LinkA{aes()}{aes()}}. If specified and
\code{inherit.aes = TRUE} (the default), it is combined with the default mapping
at the top level of the plot. You must supply \code{mapping} if there is no plot
mapping.

\item[\code{data}] The data to be displayed in this layer. There are three
options:

If \code{NULL}, the default, the data is inherited from the plot
data as specified in the call to \code{\LinkA{ggplot()}{ggplot()}}.

A \code{data.frame}, or other object, will override the plot
data. All objects will be fortified to produce a data frame. See
\code{\LinkA{fortify()}{fortify()}} for which variables will be created.

A \code{function} will be called with a single argument,
the plot data. The return value must be a \code{data.frame}, and
will be used as the layer data. A \code{function} can be created
from a \code{formula} (e.g. \code{\textasciitilde{} head(.x, 10)}).

\item[\code{stat}] The statistical transformation to use on the data for this
layer, either as a \code{ggproto} \code{Geom} subclass or as a string naming the
stat stripped of the \code{stat\_} prefix (e.g. \code{"count"} rather than
\code{"stat\_count"})

\item[\code{position}] Position adjustment, either as a string, or the result of
a call to a position adjustment function. Cannot be jointy specified with
\code{nudge\_x} or \code{nudge\_y}.

\item[\code{...}] Other arguments passed on to \code{\LinkA{layer()}{layer()}}. These are
often aesthetics, used to set an aesthetic to a fixed value, like
\code{colour = "red"} or \code{size = 3}. They may also be parameters
to the paired geom/stat.

\item[\code{parse}] If \code{TRUE}, the labels will be parsed into expressions and
displayed as described in \code{?plotmath}.

\item[\code{nudge\_x}, \code{nudge\_y}] Horizontal and vertical adjustment to nudge l
abels by. Useful for offsetting text from points, particularly
on discrete scales.

\item[\code{check\_overlap}] If \code{TRUE}, text that overlaps previous text in the
same layer will not be plotted. \code{check\_overlap} happens at draw time and in
the order of the data. Therefore data should be arranged by the label
column before calling \code{geom\_text()}. Note that this argument is not
supported by \code{geom\_label()}.

\item[\code{na.rm}] If \code{FALSE}, the default, missing values are removed with
a warning. If \code{TRUE}, missing values are silently removed.

\item[\code{show.legend}] logical. Should this layer be included in the legends?
\code{NA}, the default, includes if any aesthetics are mapped.
\code{FALSE} never includes, and \code{TRUE} always includes.
It can also be a named logical vector to finely select the aesthetics to
display.

\item[\code{inherit.aes}] If \code{FALSE}, overrides the default aesthetics,
rather than combining with them. This is most useful for helper functions
that define both data and aesthetics and shouldn't inherit behaviour from
the default plot specification, e.g. \code{\LinkA{borders()}{borders()}}.
\end{ldescription}
\end{Arguments}
%
\begin{Details}
The package will also take care of loading the StateFace font for
PDF and other devices, but to use it with the on-screen ggplot2
device, you'll need to install the font on your system.

\code{ggalt} ships with a copy of the StateFace TTF font. You can
run \code{show\_stateface()} to get the filesystem location and then
load the font manually from there.




A sample of the output from \code{geom\_stateface()}:

\Figure{geomstateface01.png}{width=10cm}

\end{Details}
%
\begin{SeeAlso}
Other StateFace operations: 
\code{\LinkA{load\_stateface}{load.Rul.stateface}()},
\code{\LinkA{show\_stateface}{show.Rul.stateface}()}
\end{SeeAlso}
%
\begin{Examples}
\begin{ExampleCode}
## Not run: 
library(ggplot2)
library(ggalt)

# Run show_stateface() to see the location of the TTF StateFace font
# You need to install it for it to work

set.seed(1492)
dat <- data.frame(state=state.abb,
                  x=sample(100, 50),
                  y=sample(100, 50),
                  col=sample(c("#b2182b", "#2166ac"), 50, replace=TRUE),
                  sz=sample(6:15, 50, replace=TRUE),
                  stringsAsFactors=FALSE)
gg <- ggplot(dat, aes(x=x, y=y))
gg <- gg + geom_stateface(aes(label=state, color=col, size=sz))
gg <- gg + scale_color_identity()
gg <- gg + scale_size_identity()
gg

## End(Not run)
\end{ExampleCode}
\end{Examples}
\HeaderA{geom\_ubar}{Uniform "bar" charts}{geom.Rul.ubar}
%
\begin{Description}
I've been using \code{geom\_segment} more to make "bar" charts, setting
\code{xend} to whatever \code{x} is and \code{yend} to \code{0}. The bar widths remain
constant without any tricks and you have granular control over the
segment width. I decided it was time to make a \code{geom}.
\end{Description}
%
\begin{Usage}
\begin{verbatim}
geom_ubar(
  mapping = NULL,
  data = NULL,
  stat = "identity",
  position = "identity",
  ...,
  na.rm = FALSE,
  show.legend = NA,
  inherit.aes = TRUE
)
\end{verbatim}
\end{Usage}
%
\begin{Arguments}
\begin{ldescription}
\item[\code{mapping}] Set of aesthetic mappings created by \code{\LinkA{aes()}{aes()}}. If specified and
\code{inherit.aes = TRUE} (the default), it is combined with the default mapping
at the top level of the plot. You must supply \code{mapping} if there is no plot
mapping.

\item[\code{data}] The data to be displayed in this layer. There are three
options:

If \code{NULL}, the default, the data is inherited from the plot
data as specified in the call to \code{\LinkA{ggplot()}{ggplot()}}.

A \code{data.frame}, or other object, will override the plot
data. All objects will be fortified to produce a data frame. See
\code{\LinkA{fortify()}{fortify()}} for which variables will be created.

A \code{function} will be called with a single argument,
the plot data. The return value must be a \code{data.frame}, and
will be used as the layer data. A \code{function} can be created
from a \code{formula} (e.g. \code{\textasciitilde{} head(.x, 10)}).

\item[\code{stat}] The statistical transformation to use on the data for this
layer, either as a \code{ggproto} \code{Geom} subclass or as a string naming the
stat stripped of the \code{stat\_} prefix (e.g. \code{"count"} rather than
\code{"stat\_count"})

\item[\code{position}] Position adjustment, either as a string naming the adjustment
(e.g. \code{"jitter"} to use \code{position\_jitter}), or the result of a call to a
position adjustment function. Use the latter if you need to change the
settings of the adjustment.

\item[\code{...}] other arguments passed on to \code{layer}. These are
often aesthetics, used to set an aesthetic to a fixed value, like
\code{color = "red"} or \code{size = 3}. They may also be parameters
to the paired geom/stat.

\item[\code{na.rm}] If \code{FALSE} (the default), removes missing values with
a warning.  If \code{TRUE} silently removes missing values.

\item[\code{show.legend}] logical. Should this layer be included in the legends?
\code{NA}, the default, includes if any aesthetics are mapped.
\code{FALSE} never includes, and \code{TRUE} always includes.
It can also be a named logical vector to finely select the aesthetics to
display.

\item[\code{inherit.aes}] If \code{FALSE}, overrides the default aesthetics,
rather than combining with them. This is most useful for helper functions
that define both data and aesthetics and shouldn't inherit behaviour from
the default plot specification, e.g. \code{\LinkA{borders()}{borders()}}.
\end{ldescription}
\end{Arguments}
%
\begin{Section}{Aesthetics}

`geom\_ubar`` understands the following aesthetics (required aesthetics are in bold):
\begin{itemize}

\item{} \strong{\code{x}}
\item{} \strong{\code{y}}
\item{} \code{alpha}
\item{} \code{colour}
\item{} \code{group}
\item{} \code{linetype}
\item{} \code{size}

\end{itemize}

\end{Section}
%
\begin{Examples}
\begin{ExampleCode}
library(ggplot2)

data(economics)
ggplot(economics, aes(date, uempmed)) +
  geom_ubar()
\end{ExampleCode}
\end{Examples}
\HeaderA{geom\_xspline}{Connect control points/observations with an X-spline}{geom.Rul.xspline}
\aliasA{stat\_xspline}{geom\_xspline}{stat.Rul.xspline}
\keyword{xspline implementations}{geom\_xspline}
%
\begin{Description}
Draw an X-spline, a curve drawn relative to control points/observations.
Patterned after \code{geom\_line} in that it orders the points by \code{x}
first before computing the splines.
\end{Description}
%
\begin{Usage}
\begin{verbatim}
geom_xspline(
  mapping = NULL,
  data = NULL,
  stat = "xspline",
  position = "identity",
  na.rm = TRUE,
  show.legend = NA,
  inherit.aes = TRUE,
  spline_shape = -0.25,
  open = TRUE,
  rep_ends = TRUE,
  ...
)

stat_xspline(
  mapping = NULL,
  data = NULL,
  geom = "line",
  position = "identity",
  na.rm = TRUE,
  show.legend = NA,
  inherit.aes = TRUE,
  spline_shape = -0.25,
  open = TRUE,
  rep_ends = TRUE,
  ...
)
\end{verbatim}
\end{Usage}
%
\begin{Arguments}
\begin{ldescription}
\item[\code{mapping}] Set of aesthetic mappings created by \code{\LinkA{aes()}{aes()}}. If specified and
\code{inherit.aes = TRUE} (the default), it is combined with the default mapping
at the top level of the plot. You must supply \code{mapping} if there is no plot
mapping.

\item[\code{data}] The data to be displayed in this layer. There are three
options:

If \code{NULL}, the default, the data is inherited from the plot
data as specified in the call to \code{\LinkA{ggplot()}{ggplot()}}.

A \code{data.frame}, or other object, will override the plot
data. All objects will be fortified to produce a data frame. See
\code{\LinkA{fortify()}{fortify()}} for which variables will be created.

A \code{function} will be called with a single argument,
the plot data. The return value must be a \code{data.frame}, and
will be used as the layer data. A \code{function} can be created
from a \code{formula} (e.g. \code{\textasciitilde{} head(.x, 10)}).

\item[\code{position}] Position adjustment, either as a string naming the adjustment
(e.g. \code{"jitter"} to use \code{position\_jitter}), or the result of a call to a
position adjustment function. Use the latter if you need to change the
settings of the adjustment.

\item[\code{na.rm}] If \code{FALSE}, the default, missing values are removed with
a warning. If \code{TRUE}, missing values are silently removed.

\item[\code{show.legend}] logical. Should this layer be included in the legends?
\code{NA}, the default, includes if any aesthetics are mapped.
\code{FALSE} never includes, and \code{TRUE} always includes.
It can also be a named logical vector to finely select the aesthetics to
display.

\item[\code{inherit.aes}] If \code{FALSE}, overrides the default aesthetics,
rather than combining with them. This is most useful for helper functions
that define both data and aesthetics and shouldn't inherit behaviour from
the default plot specification, e.g. \code{\LinkA{borders()}{borders()}}.

\item[\code{spline\_shape}] A numeric vector of values between -1 and 1, which
control the shape of the spline relative to the control points.

\item[\code{open}] A logical value indicating whether the spline is an open or a
closed shape.

\item[\code{rep\_ends}] For open X-splines, a logical value indicating whether the
first and last control points should be replicated for drawing the
curve. Ignored for closed X-splines.

\item[\code{...}] Other arguments passed on to \code{\LinkA{layer()}{layer()}}. These are
often aesthetics, used to set an aesthetic to a fixed value, like
\code{colour = "red"} or \code{size = 3}. They may also be parameters
to the paired geom/stat.

\item[\code{geom}, \code{stat}] Use to override the default connection between
\code{geom\_xspline} and \code{stat\_xspline}.
\end{ldescription}
\end{Arguments}
%
\begin{Details}



A sample of the output from \code{geom\_xspline()}:

\Figure{geomxspline01.png}{width=10cm}



An X-spline is a line drawn relative to control points. For each control
point, the line may pass through (interpolate) the control point or it may
only approach (approximate) the control point; the behaviour is determined
by a shape parameter for each control point.

If the shape parameter is greater than zero, the spline approximates the
control points (and is very similar to a cubic B-spline when the shape is
1). If the shape parameter is less than zero, the spline interpolates the
control points (and is very similar to a Catmull-Rom spline when the shape
is -1). If the shape parameter is 0, the spline forms a sharp corner at that
control point.

For open X-splines, the start and end control points must have a shape of
0 (and non-zero values are silently converted to zero).

For open X-splines, by default the start and end control points are
replicated before the curve is drawn. A curve is drawn between (interpolating
or approximating) the second and third of each set of four control points,
so this default behaviour ensures that the resulting curve starts at the
first control point you have specified and ends at the last control point.
The default behaviour can be turned off via the repEnds argument.
\end{Details}
%
\begin{Section}{Aesthetics}

\code{geom\_xspline} understands the following aesthetics (required aesthetics
are in bold):
\begin{itemize}

\item{} \strong{\code{x}}
\item{} \strong{\code{y}}
\item{} \code{alpha}
\item{} \code{color}
\item{} \code{linetype}
\item{} \code{size}

\end{itemize}

\end{Section}
%
\begin{Section}{Computed variables}

\begin{itemize}

\item{} x
\item{} y

\end{itemize}

\end{Section}
%
\begin{References}
Blanc, C. and Schlick, C. (1995), "X-splines : A Spline Model
Designed for the End User", in \emph{Proceedings of SIGGRAPH 95},
pp. 377-386. \url{http://dept-info.labri.fr/~schlick/DOC/sig1.html}
\end{References}
%
\begin{SeeAlso}
\code{\LinkA{geom\_line}{geom.Rul.line}}: Connect observations (x order);
\code{\LinkA{geom\_path}{geom.Rul.path}}: Connect observations;
\code{\LinkA{geom\_polygon}{geom.Rul.polygon}}: Filled paths (polygons);
\code{\LinkA{geom\_segment}{geom.Rul.segment}}: Line segments;
\code{\LinkA{xspline}{xspline}};
\code{\LinkA{grid.xspline}{grid.xspline}}

Other xspline implementations: 
\code{\LinkA{geom\_xspline2}{geom.Rul.xspline2}()}
\end{SeeAlso}
%
\begin{Examples}
\begin{ExampleCode}
set.seed(1492)
dat <- data.frame(x=c(1:10, 1:10, 1:10),
                  y=c(sample(15:30, 10), 2*sample(15:30, 10),
                      3*sample(15:30, 10)),
                  group=factor(c(rep(1, 10), rep(2, 10), rep(3, 10)))
)

ggplot(dat, aes(x, y, group=group, color=group)) +
  geom_point() +
  geom_line()

ggplot(dat, aes(x, y, group=group, color=factor(group))) +
  geom_point() +
  geom_line() +
  geom_smooth(se=FALSE, linetype="dashed", size=0.5)

ggplot(dat, aes(x, y, group=group, color=factor(group))) +
  geom_point(color="black") +
  geom_smooth(se=FALSE, linetype="dashed", size=0.5) +
  geom_xspline(size=0.5)

ggplot(dat, aes(x, y, group=group, color=factor(group))) +
  geom_point(color="black") +
  geom_smooth(se=FALSE, linetype="dashed", size=0.5) +
  geom_xspline(spline_shape=-0.4, size=0.5)

ggplot(dat, aes(x, y, group=group, color=factor(group))) +
  geom_point(color="black") +
  geom_smooth(se=FALSE, linetype="dashed", size=0.5) +
  geom_xspline(spline_shape=0.4, size=0.5)

ggplot(dat, aes(x, y, group=group, color=factor(group))) +
  geom_point(color="black") +
  geom_smooth(se=FALSE, linetype="dashed", size=0.5) +
  geom_xspline(spline_shape=1, size=0.5)

ggplot(dat, aes(x, y, group=group, color=factor(group))) +
  geom_point(color="black") +
  geom_smooth(se=FALSE, linetype="dashed", size=0.5) +
  geom_xspline(spline_shape=0, size=0.5)

ggplot(dat, aes(x, y, group=group, color=factor(group))) +
  geom_point(color="black") +
  geom_smooth(se=FALSE, linetype="dashed", size=0.5) +
  geom_xspline(spline_shape=-1, size=0.5)
\end{ExampleCode}
\end{Examples}
\HeaderA{geom\_xspline2}{Alternative implemenation for connecting  control points/observations with an X-spline}{geom.Rul.xspline2}
\keyword{xspline implementations}{geom\_xspline2}
%
\begin{Description}
Alternative implemenation for connecting  control points/observations
with an X-spline
\end{Description}
%
\begin{Usage}
\begin{verbatim}
geom_xspline2(
  mapping = NULL,
  data = NULL,
  stat = "identity",
  position = "identity",
  na.rm = FALSE,
  show.legend = NA,
  inherit.aes = TRUE,
  ...
)
\end{verbatim}
\end{Usage}
%
\begin{Arguments}
\begin{ldescription}
\item[\code{mapping}] Set of aesthetic mappings created by \code{\LinkA{aes()}{aes()}}. If specified and
\code{inherit.aes = TRUE} (the default), it is combined with the default mapping
at the top level of the plot. You must supply \code{mapping} if there is no plot
mapping.

\item[\code{data}] The data to be displayed in this layer. There are three
options:

If \code{NULL}, the default, the data is inherited from the plot
data as specified in the call to \code{\LinkA{ggplot()}{ggplot()}}.

A \code{data.frame}, or other object, will override the plot
data. All objects will be fortified to produce a data frame. See
\code{\LinkA{fortify()}{fortify()}} for which variables will be created.

A \code{function} will be called with a single argument,
the plot data. The return value must be a \code{data.frame}, and
will be used as the layer data. A \code{function} can be created
from a \code{formula} (e.g. \code{\textasciitilde{} head(.x, 10)}).

\item[\code{position}] Position adjustment, either as a string naming the adjustment
(e.g. \code{"jitter"} to use \code{position\_jitter}), or the result of a call to a
position adjustment function. Use the latter if you need to change the
settings of the adjustment.

\item[\code{na.rm}] If \code{FALSE}, the default, missing values are removed with
a warning. If \code{TRUE}, missing values are silently removed.

\item[\code{show.legend}] logical. Should this layer be included in the legends?
\code{NA}, the default, includes if any aesthetics are mapped.
\code{FALSE} never includes, and \code{TRUE} always includes.
It can also be a named logical vector to finely select the aesthetics to
display.

\item[\code{inherit.aes}] If \code{FALSE}, overrides the default aesthetics,
rather than combining with them. This is most useful for helper functions
that define both data and aesthetics and shouldn't inherit behaviour from
the default plot specification, e.g. \code{\LinkA{borders()}{borders()}}.

\item[\code{...}] Other arguments passed on to \code{\LinkA{layer()}{layer()}}. These are
often aesthetics, used to set an aesthetic to a fixed value, like
\code{colour = "red"} or \code{size = 3}. They may also be parameters
to the paired geom/stat.
\end{ldescription}
\end{Arguments}
%
\begin{Value}
creates a spline curve
\end{Value}
%
\begin{Author}
Ben Bolker
\end{Author}
%
\begin{SeeAlso}
Other xspline implementations: 
\code{\LinkA{geom\_xspline}{geom.Rul.xspline}()}
\end{SeeAlso}
\HeaderA{ggalt}{Extra Geoms, Stats, Coords, Scales \& Fonts for 'ggplot2'}{ggalt}
%
\begin{Description}
A package containing additional geoms, coords, stats, scales \& fonts
for ggplot2 2.0+
\end{Description}
%
\begin{Author}
Bob Rudis (@hrbrmstr)
\end{Author}
\HeaderA{load\_stateface}{Load stateface font}{load.Rul.stateface}
\keyword{StateFace operations}{load\_stateface}
%
\begin{Description}
Makes the ProPublica StateFace font available to PDF, PostScript,
et. al. devices.
\end{Description}
%
\begin{Usage}
\begin{verbatim}
load_stateface()
\end{verbatim}
\end{Usage}
%
\begin{SeeAlso}
Other StateFace operations: 
\code{\LinkA{geom\_stateface}{geom.Rul.stateface}()},
\code{\LinkA{show\_stateface}{show.Rul.stateface}()}
\end{SeeAlso}
\HeaderA{plotly\_helpers}{Plotly helpers}{plotly.Rul.helpers}
\aliasA{to\_basic.GeomBkde2d}{plotly\_helpers}{to.Rul.basic.GeomBkde2d}
\aliasA{to\_basic.GeomStateface}{plotly\_helpers}{to.Rul.basic.GeomStateface}
\aliasA{to\_basic.GeomXspline}{plotly\_helpers}{to.Rul.basic.GeomXspline}
\keyword{internal}{plotly\_helpers}
%
\begin{Description}
Helper functions to make it easier to automatically create plotly charts
\end{Description}
%
\begin{Usage}
\begin{verbatim}
to_basic.GeomXspline(data, prestats_data, layout, params, p, ...)

to_basic.GeomBkde2d(data, prestats_data, layout, params, p, ...)

to_basic.GeomStateface(data, prestats_data, layout, params, p, ...)
\end{verbatim}
\end{Usage}
%
\begin{Arguments}
\begin{ldescription}
\item[\code{data}, \code{prestats\_data}, \code{layout}, \code{params}, \code{p}, \code{...}] plotly interface parameters
\end{ldescription}
\end{Arguments}
\HeaderA{position\_dodgev}{Vertically dodge position}{position.Rul.dodgev}
\aliasA{PositionDodgev}{position\_dodgev}{PositionDodgev}
\keyword{datasets}{position\_dodgev}
%
\begin{Description}
Vertically dodge position
\end{Description}
%
\begin{Usage}
\begin{verbatim}
position_dodgev(height = NULL)
\end{verbatim}
\end{Usage}
%
\begin{Arguments}
\begin{ldescription}
\item[\code{height}] numeric, height of vertical dodge, Default: NULL
\end{ldescription}
\end{Arguments}
%
\begin{Note}
position-dodgev(): unmodified from lionel-/ggstance/R/position-dodgev.R 73f521384ae8ea277db5f7d5a2854004aa18f947
\end{Note}
%
\begin{Author}
@ggstance authors
\end{Author}
%
\begin{Examples}
\begin{ExampleCode}

if(interactive()){

dat <- data.frame(
  trt = c(LETTERS[1:5], "D"),
  l = c(20, 40, 10, 30, 50, 40),
  r = c(70, 50, 30, 60, 80, 70)
)

ggplot(dat, aes(y=trt, x=l, xend=r)) +
 geom_dumbbell(size=3, color="#e3e2e1",
               colour_x = "#5b8124", colour_xend = "#bad744",
               dot_guide=TRUE, dot_guide_size=0.25,
               position=position_dodgev(height=0.8)) +
 labs(x=NULL, y=NULL, title="ggplot2 geom_dumbbell with dot guide") +
 theme_minimal() +
 theme(panel.grid.major.x=element_line(size=0.05))

 }

\end{ExampleCode}
\end{Examples}
\HeaderA{show\_stateface}{Show location of StateFace font}{show.Rul.stateface}
\keyword{StateFace operations}{show\_stateface}
%
\begin{Description}
Displays the path to the StateFace font. For the font to work
in the on-screen plot device for ggplot2, you need to install
the font on your system
\end{Description}
%
\begin{Usage}
\begin{verbatim}
show_stateface()
\end{verbatim}
\end{Usage}
%
\begin{SeeAlso}
Other StateFace operations: 
\code{\LinkA{geom\_stateface}{geom.Rul.stateface}()},
\code{\LinkA{load\_stateface}{load.Rul.stateface}()}
\end{SeeAlso}
\HeaderA{stat\_ash}{Compute and display a univariate averaged shifted histogram (polynomial kernel)}{stat.Rul.ash}
%
\begin{Description}
See \code{\LinkA{bin1}{bin1}} \& \code{\LinkA{ash1}{ash1}} for more information.
\end{Description}
%
\begin{Usage}
\begin{verbatim}
stat_ash(
  mapping = NULL,
  data = NULL,
  geom = "area",
  position = "stack",
  ab = NULL,
  nbin = 50,
  m = 5,
  kopt = c(2, 2),
  na.rm = FALSE,
  show.legend = NA,
  inherit.aes = TRUE,
  ...
)
\end{verbatim}
\end{Usage}
%
\begin{Arguments}
\begin{ldescription}
\item[\code{mapping}] Set of aesthetic mappings created by \code{\LinkA{aes()}{aes()}}. If specified and
\code{inherit.aes = TRUE} (the default), it is combined with the default mapping
at the top level of the plot. You must supply \code{mapping} if there is no plot
mapping.

\item[\code{data}] The data to be displayed in this layer. There are three
options:

If \code{NULL}, the default, the data is inherited from the plot
data as specified in the call to \code{\LinkA{ggplot()}{ggplot()}}.

A \code{data.frame}, or other object, will override the plot
data. All objects will be fortified to produce a data frame. See
\code{\LinkA{fortify()}{fortify()}} for which variables will be created.

A \code{function} will be called with a single argument,
the plot data. The return value must be a \code{data.frame}, and
will be used as the layer data. A \code{function} can be created
from a \code{formula} (e.g. \code{\textasciitilde{} head(.x, 10)}).

\item[\code{geom}] Use to override the default Geom

\item[\code{position}] Position adjustment, either as a string naming the adjustment
(e.g. \code{"jitter"} to use \code{position\_jitter}), or the result of a call to a
position adjustment function. Use the latter if you need to change the
settings of the adjustment.

\item[\code{ab}] half-open interval for bins \emph{[a,b)}. If no value is specified,
the range of x is stretched by \code{5\%} at each end and used the
interval.

\item[\code{nbin}] number of bins desired. Default \code{50}.

\item[\code{m}] integer smoothing parameter; Default \code{5}.

\item[\code{kopt}] vector of length 2 specifying the kernel, which is proportional
to \emph{( 1 - abs(i/m)\textasciicircum{}kopt(1) )i\textasciicircum{}kopt(2)}; (2,2)=biweight (default);
(0,0)=uniform; (1,0)=triangle; (2,1)=Epanechnikov; (2,3)=triweight.

\item[\code{na.rm}] If \code{FALSE}, the default, missing values are removed with
a warning. If \code{TRUE}, missing values are silently removed.

\item[\code{show.legend}] logical. Should this layer be included in the legends?
\code{NA}, the default, includes if any aesthetics are mapped.
\code{FALSE} never includes, and \code{TRUE} always includes.
It can also be a named logical vector to finely select the aesthetics to
display.

\item[\code{inherit.aes}] If \code{FALSE}, overrides the default aesthetics,
rather than combining with them. This is most useful for helper functions
that define both data and aesthetics and shouldn't inherit behaviour from
the default plot specification, e.g. \code{\LinkA{borders()}{borders()}}.

\item[\code{...}] Other arguments passed on to \code{\LinkA{layer()}{layer()}}. These are
often aesthetics, used to set an aesthetic to a fixed value, like
\code{colour = "red"} or \code{size = 3}. They may also be parameters
to the paired geom/stat.
\end{ldescription}
\end{Arguments}
%
\begin{Details}



A sample of the output from \code{stat\_ash()}:

\Figure{statash01.png}{width=10cm}

\end{Details}
%
\begin{Section}{Aesthetics}

\code{geom\_ash} understands the following aesthetics (required aesthetics
are in bold):
\begin{itemize}

\item{} \strong{\code{x}}
\item{} \code{alpha}
\item{} \code{color}
\item{} \code{fill}
\item{} \code{linetype}
\item{} \code{size}

\end{itemize}

\end{Section}
%
\begin{Section}{Computed variables}

\begin{description}

\item[\code{density}] ash density estimate

\end{description}

\end{Section}
%
\begin{References}
David Scott (1992), \emph{"Multivariate Density Estimation,"}
John Wiley, (chapter 5 in particular).\\{}
\\{}
B. W. Silverman (1986), \emph{"Density Estimation for Statistics
and Data Analysis,"} Chapman \& Hall.
\end{References}
%
\begin{Examples}
\begin{ExampleCode}
# compare
library(gridExtra)
set.seed(1492)
dat <- data.frame(x=rnorm(100))
grid.arrange(ggplot(dat, aes(x)) + stat_ash(),
             ggplot(dat, aes(x)) + stat_bkde(),
             ggplot(dat, aes(x)) + stat_density(),
             nrow=3)

cols <- RColorBrewer::brewer.pal(3, "Dark2")
ggplot(dat, aes(x)) +
  stat_ash(alpha=1/2, fill=cols[3]) +
  stat_bkde(alpha=1/2, fill=cols[2]) +
  stat_density(alpha=1/2, fill=cols[1]) +
  geom_rug() +
  labs(x=NULL, y="density/estimate") +
  scale_x_continuous(expand=c(0,0)) +
  theme_bw() +
  theme(panel.grid=element_blank()) +
  theme(panel.border=element_blank())
\end{ExampleCode}
\end{Examples}
\HeaderA{stat\_stepribbon}{Step ribbon statistic}{stat.Rul.stepribbon}
%
\begin{Description}
Provides stairstep values for ribbon plots
\end{Description}
%
\begin{Usage}
\begin{verbatim}
stat_stepribbon(
  mapping = NULL,
  data = NULL,
  geom = "ribbon",
  position = "identity",
  na.rm = FALSE,
  show.legend = NA,
  inherit.aes = TRUE,
  direction = "hv",
  ...
)
\end{verbatim}
\end{Usage}
%
\begin{Arguments}
\begin{ldescription}
\item[\code{mapping}] Set of aesthetic mappings created by \code{\LinkA{aes()}{aes()}}. If specified and
\code{inherit.aes = TRUE} (the default), it is combined with the default mapping
at the top level of the plot. You must supply \code{mapping} if there is no plot
mapping.

\item[\code{data}] The data to be displayed in this layer. There are three
options:

If \code{NULL}, the default, the data is inherited from the plot
data as specified in the call to \code{\LinkA{ggplot()}{ggplot()}}.

A \code{data.frame}, or other object, will override the plot
data. All objects will be fortified to produce a data frame. See
\code{\LinkA{fortify()}{fortify()}} for which variables will be created.

A \code{function} will be called with a single argument,
the plot data. The return value must be a \code{data.frame}, and
will be used as the layer data. A \code{function} can be created
from a \code{formula} (e.g. \code{\textasciitilde{} head(.x, 10)}).

\item[\code{geom}] which geom to use; defaults to "\code{ribbon}"

\item[\code{position}] Position adjustment, either as a string naming the adjustment
(e.g. \code{"jitter"} to use \code{position\_jitter}), or the result of a call to a
position adjustment function. Use the latter if you need to change the
settings of the adjustment.

\item[\code{na.rm}] If \code{FALSE}, the default, missing values are removed with
a warning. If \code{TRUE}, missing values are silently removed.

\item[\code{show.legend}] logical. Should this layer be included in the legends?
\code{NA}, the default, includes if any aesthetics are mapped.
\code{FALSE} never includes, and \code{TRUE} always includes.
It can also be a named logical vector to finely select the aesthetics to
display.

\item[\code{inherit.aes}] If \code{FALSE}, overrides the default aesthetics,
rather than combining with them. This is most useful for helper functions
that define both data and aesthetics and shouldn't inherit behaviour from
the default plot specification, e.g. \code{\LinkA{borders()}{borders()}}.

\item[\code{direction}] \code{hv} for horizontal-veritcal steps, \code{vh} for
vertical-horizontal steps

\item[\code{...}] Other arguments passed on to \code{\LinkA{layer()}{layer()}}. These are
often aesthetics, used to set an aesthetic to a fixed value, like
\code{colour = "red"} or \code{size = 3}. They may also be parameters
to the paired geom/stat.
\end{ldescription}
\end{Arguments}
%
\begin{References}
\url{https://groups.google.com/forum/?fromgroups=\#!topic/ggplot2/9cFWHaH1CPs}
\end{References}
%
\begin{Examples}
\begin{ExampleCode}
x <- 1:10
df <- data.frame(x=x, y=x+10, ymin=x+7, ymax=x+12)

gg <- ggplot(df, aes(x, y))
gg <- gg + geom_ribbon(aes(ymin=ymin, ymax=ymax),
                       stat="stepribbon", fill="#b2b2b2")
gg <- gg + geom_step(color="#2b2b2b")
gg

gg <- ggplot(df, aes(x, y))
gg <- gg + geom_ribbon(aes(ymin=ymin, ymax=ymax),
                       stat="stepribbon", fill="#b2b2b2",
                       direction="hv")
gg <- gg + geom_step(color="#2b2b2b")
gg
\end{ExampleCode}
\end{Examples}
\printindex{}
\end{document}
